\section{Introduction: a quiet question in the committee room}
\label{sec:intro}

This review is written for a reader who shapes medical and artificial
intelligence law but who has never driven a surgical robot or written a line of
verification code. It therefore defines its terms as it goes, and it asks one
quiet question in eight different ways: how did we ever ask human beings, alone,
to carry a task this large without help?

A few plain definitions make the rest of the story readable. \textit{Physical AI}
is artificial intelligence that acts in the physical world through a machine, such
as a surgical humanoid in an oncology trial, rather than only returning text on a
screen. \textit{Verification before generation} is the principle at the center of
H. R. 9510 \cite{Kawchak2026HR9510Bill}: a machine may propose an action, but the
proposal is checked against fixed, published standards before anything reaches a
patient. \textit{Claude Code} is the software agent that drafts the candidate
procedure, \textit{Codex} is a second, independent agent that reviews it, and
\textit{VVUQ}, meaning verification, validation, and uncertainty quantification,
is the ten-gate examination the proposal must pass. The examination resolves in
exactly one of three ways: ACCEPT and execute under audit, ESCALATE to a qualified
human, or BLOCK before execution. Figure 1 shows the whole
mechanism at a glance.

 
\begin{narrativefig}
\node[nfone] (a) at (0,0) {Human\\specification};
\node[nftwo] (b) at (4,0) {Claude Code\\generation};
\node[nftwo] (c) at (8,0) {Codex\\review};
\node[nfdec] (d) at (12,0) {VVUQ ten gates};
\node[nfhope] (acc) at (5.0,-2.6) {ACCEPT under audit};
\node[nfthree] (esc) at (8.4,-2.6) {ESCALATE to human};
\node[nfrisk] (blk) at (11.8,-2.6) {BLOCK execution};
\draw[nfedge] (a)--(b); \draw[nfedge] (b)--(c); \draw[nfedge] (c)--(d);
\draw[nfedge] (d) to[out=-160,in=60, looseness=0.4] (acc.north);
\draw[nfedge] (d) -- (esc.north);
\draw[nfedge] (d) to[out=-60,in=70, looseness=0.2] (blk.north);
\draw[nfedged] (esc.north) to[out=0,in=-90,looseness=0.8] (b.south);
\end{narrativefig}
\vspace{-0.65cm}
\figcaption{Figure 1. Verification before generation. A human specification is drafted by one agent, reviewed by a
second, and admitted only after ten gates resolve to ACCEPT, ESCALATE, or BLOCK.}

The three guiding questions return throughout. How did we ever utilize humans
alone for this challenging task? Why did we ever fully trust the prior human
clinical-trial system if there were so many opportunities for compounding human
error? And why would the human peer-review process continue to be used for
increasingly complex and voluminous artificial intelligence outputs? These are
asked politely, not to diminish the clinicians who have carried the system on
their shoulders, but to notice that the load has outgrown unaided human attention.

The advocacy literature is clear that emotion and evidence must travel together.
Legislators are moved by compelling, concise stories that make an abstract
technology tangible \cite{novak2020patient}, yet they weigh the credibility of the
presenter and the underlying data at least as heavily as the intensity of the
appeal \cite{moreland2015hearing}. This review therefore organizes its case as
eight emotional pillars, and pairs each with a cited fact.
Figure 2 previews that path, from the most common appeal,
compassion, to the closing appeal, urgency, and to single action.

\begin{narrativefig}
\node[nfone,text width=16mm]  (p1) at (0,0)    {1\\Compassion};
\node[nfrisk,text width=16mm] (p2) at (3.0,0)  {2\\Fear};
\node[nfrisk,text width=16mm] (p3) at (6.0,0)  {3\\Outrage};
\node[nfhope,text width=16mm] (p4) at (9.0,0)  {4\\Hope};
\node[nftwo,text width=16mm]  (p5) at (9.0,-1.8){5\\Duty};
\node[nftwo,text width=16mm]  (p6) at (6.0,-1.8){6\\Protect};
\node[nfhope,text width=16mm] (p7) at (3.0,-1.8){7\\Trust};
\node[nfrisk,text width=16mm] (p8) at (0,-1.8) {8\\Urgency};
\node[nfact,text width=22mm]  (ax) at (4.5,-3.7){Enact H. R. 9510};
\draw[nfedge] (p1)--(p2); \draw[nfedge] (p2)--(p3); \draw[nfedge] (p3)--(p4);
\draw[nfedge] (p4)--(p5); \draw[nfedge] (p5)--(p6); \draw[nfedge] (p6)--(p7);
\draw[nfedge] (p7)--(p8); \draw[nfedge] (p8) to[out=-90,in=180] (ax.west);
\end{narrativefig}
\vspace{-0.75cm}
\figcaption{Figure 2. The eight emotional pillars. The argument accumulates from compassion to urgency and converges
on a single legislative request.}

The remainder of the review takes the pillars in turn. Each section opens with a
human appeal, anchors it in credible evidence, defines any term on first use, and
closes by naming the specific provision of H. R. 9510 that answers it.
